\chapter{Il Piedistallo Cognitivo: Quando l'Ammirazione Acceca}

\begin{quote}
	\textit{``Il problema con il mondo è che gli stupidi sono sicuri di sé mentre gli intelligenti sono pieni di dubbi. E noi tendiamo a credere a chi sembra più sicuro.''}\\
	\hfill Bertrand Russell (parafrasato)
\end{quote}

Nel 1954, Linus Pauling era l'uomo più decorato della scienza americana: due premi Nobel, un curriculum che faceva impallidire chiunque, e una reputazione di infallibilità. Quando iniziò a sostenere che dosi massicce di vitamina C potessero curare il raffreddore, il cancro e praticamente qualsiasi malattia, il mondo scientifico si trovò di fronte a un dilemma imbarazzante. I dati non lo sostenevano. Le sperimentazioni lo smentivano. Eppure, per anni, pochissimi ebbero il coraggio di dire pubblicamente quello che tutti sapevano in privato: Pauling, su questo punto, aveva torto marcio.\footnote{Pauling, L. (1970). \textit{Vitamin C and the Common Cold}. W.H. Freeman. Le meta-analisi successive hanno smentito le sue affermazioni sugli effetti terapeutici delle megadosi di vitamina C.}

Perché? Cosa impedisce a un'intera comunità di esperti di contraddire un collega quando sbaglia? La risposta non sta nella cortesia professionale, né nella paura delle ritorsioni accademiche. Sta in qualcosa di molto più profondo: un cortocircuito del nostro cervello, un meccanismo che, quando si attiva, rende quasi impossibile vedere gli errori delle persone che ammiriamo.

Benvenuti sul piedistallo cognitivo.

\section{Un filtro invisibile nel cervello}

Nel 1920, lo psicologo Edward Thorndike fece un esperimento apparentemente banale: chiese a degli ufficiali dell'esercito di valutare i propri soldati su una serie di caratteristiche indipendenti, dall'intelligenza alla forma fisica, dalla leadership alla puntualità. Quello che scoprì fu sorprendente. Se un ufficiale giudicava un soldato di bell'aspetto, tendeva automaticamente a giudicarlo anche più intelligente, più disciplinato, più capace. Come se la bellezza ``irradiasse'' luce su tutte le altre qualità. Thorndike lo chiamò \textit{halo effect}, effetto alone.\footnote{Thorndike, E. L. (1920). ``A constant error in psychological ratings''. \textit{Journal of Applied Psychology}, 4(1), 25--29.}

Pensatelo come un filtro fotografico, ma per il giudizio. Quando qualcuno ci colpisce positivamente per una caratteristica (bellezza, eloquenza, successo), il nostro cervello applica quella impressione favorevole a tutto il resto, comprese aree che non abbiamo mai verificato. È un meccanismo di economia mentale: valutare ogni singola qualità di ogni persona che incontriamo richiederebbe un dispendio enorme di energia. Quindi il cervello prende una scorciatoia. ``Mi piace questa cosa di te, quindi probabilmente mi piace anche il resto.''

Fin qui, si tratta di un errore comprensibile, quasi simpatico. Il problema inizia quando questo filtro incontra un altro meccanismo, più antico e più potente: il \textit{bias di autorità}, la tendenza profondamente radicata nella nostra biologia a dare un peso sproporzionato alle parole di chi percepiamo come figura dominante o esperta.

Le neuroscienze ci spiegano perché. Quando ascoltiamo una persona che consideriamo autorevole, le aree del cervello legate alla valutazione critica, in particolare la corteccia prefrontale dorsolaterale, riducono la loro attività. È come se il cervello, di fronte a un'autorità riconosciuta, abbassasse le difese.\footnote{Milgram, S. (1963). ``Behavioral study of obedience''. \textit{Journal of Abnormal and Social Psychology}, 67(4), 371--378.} In termini evolutivi ha senso: nella tribù ancestrale, mettere in discussione il capo nel momento sbagliato poteva costare la vita. Obbedire rapidamente, invece, aumentava le probabilità di sopravvivenza.

Ma nel mondo moderno, dove le ``tribù'' sono community online e i ``capi'' sono CEO che twittano alle tre di notte, questo istinto può diventare una trappola.

\section{L'avvocato difensore che vive nella nostra testa}

Quando l'effetto alone e il bias di autorità si combinano, nasce qualcosa di particolarmente insidioso: una forma di cecità selettiva. Non si tratta semplicemente di dare più credito a chi ammiriamo. Il nostro cervello attiva un vero e proprio meccanismo di difesa dell'immagine idealizzata.

Gli psicologi cognitivi lo chiamano \textit{motivated reasoning}, ragionamento motivato. Funziona così: quando il nostro eroe intellettuale dice qualcosa di palesemente sbagliato, il cervello non registra l'errore. Si trasforma, invece, in un avvocato difensore.\footnote{Kunda, Z. (1990). ``The case for motivated reasoning''. \textit{Psychological Bulletin}, 108(3), 480--498.} ``Forse ho capito male io.'' ``Probabilmente c'è un contesto che mi sfugge.'' ``Starà semplificando per il grande pubblico.'' È come se dentro di noi si attivasse un processo in tribunale dove il nostro idolo è sempre innocente fino a prova contraria, e la prova contraria non è mai sufficiente.

Questo fenomeno è particolarmente visibile nel mondo accademico. Pensate alle conferenze scientifiche: un premio Nobel sale sul palco e presenta dati traballanti o un'ipotesi azzardata. La platea annuisce, pensierosa. Qualcuno, forse, alza un sopracciglio. Ma nessuno si alza a contestare. Non per vigliaccheria, non solo per deferenza. Ma perché il cervello di chi ascolta sta letteralmente filtrando le incongruenze, riempiendo i buchi logici con giustificazioni automatiche. Ora immaginate che la stessa identica presentazione venga fatta da un dottorando sconosciuto. Verrebbe fatta a pezzi in cinque minuti.

La differenza tra i due scenari non sta nella qualità dei dati. Sta nel piedistallo.

\section{L'illusione dell'expertise trasferibile}

C'è una credenza, tanto diffusa quanto pericolosa, secondo cui l'intelligenza sarebbe una sorta di moneta universale: se sei geniale in un campo, quella genialità si trasferisce automaticamente a qualsiasi altro ambito. Un fisico brillante deve per forza capire di economia. Un imprenditore visionario deve sapere come risolvere i problemi sociali. Un chirurgo eccezionale deve avere opinioni illuminate sulla politica estera.

Nella tribù ancestrale, questa intuizione aveva un senso. Il miglior cacciatore era probabilmente anche un buon osservatore, un individuo resistente e capace di prendere decisioni rapide. Le competenze, in un mondo semplice, si sovrapponevano davvero. Ma la modernità ha cambiato le regole del gioco. Viviamo in un mondo di specializzazione estrema, dove un cardiochirurgo di fama mondiale può non sapere nulla di virologia, e un genio della programmazione può avere idee disastrose sulla geopolitica.

Eppure, il nostro cervello continua a usare la vecchia euristica. E i risultati, a volte, sono tragici.

Steve Jobs è forse l'esempio più eloquente. Genio indiscusso del design e del marketing tecnologico, capace di intuire i desideri dei consumatori prima che loro stessi ne fossero consapevoli. Ma quando gli venne diagnosticato un tumore al pancreas, per nove mesi rifiutò l'intervento chirurgico, optando invece per diete a base di frutta, agopuntura e rimedi erboristici.\footnote{Isaacson, W. (2011). \textit{Steve Jobs}. Simon \& Schuster.} Quanti, nel suo entourage, ebbero il coraggio di dirgli chiaramente che stava commettendo un errore potenzialmente fatale? L'alone della sua genialità tecnologica si estendeva, come un campo magnetico, anche alle sue decisioni mediche. E quell'alone, in quel caso, ebbe conseguenze irreversibili.

\section{Quando il piedistallo diventa un business}

Se il piedistallo cognitivo è sempre esistito, il mondo contemporaneo lo ha trasformato in qualcosa di qualitativamente diverso. I social media, in particolare, hanno introdotto due novità che amplificano il fenomeno in modi senza precedenti.

La prima è la \textit{relazione parasociale con l'autorità}. Prima dell'era digitale, i nostri idoli erano figure distanti, mediate. Vedevamo Einstein in fotografie posate, leggevamo le sue citazioni selezionate con cura. La distanza, paradossalmente, aiutava: era più facile ammirare con lucidità qualcuno che non si aveva l'illusione di conoscere. Oggi, i nostri idoli condividono ogni pensiero che attraversa loro la mente. Li vediamo fare colazione, litigare su Twitter, pubblicare opinioni a caldo su argomenti di cui sanno poco. E invece di demistificarli, questa esposizione costante produce un effetto opposto: la sensazione di intimità. ``Non solo è un genio, ma è come se fosse un mio amico.'' Questa finta vicinanza non indebolisce l'effetto alone. Lo cementa.\footnote{Horton, D., \& Wohl, R. R. (1956). ``Mass communication and para-social interaction''. \textit{Psychiatry}, 19(3), 215--229.}

La seconda novità è ciò che potremmo chiamare \textit{l'autorità istantanea}. Qualcuno scrive un thread virale su un argomento qualsiasi (economia, epidemiologia, politica estera) e nel giro di ore diventa ``la voce autorevole'' su quel tema. Non importa se non ha credenziali, esperienza, o anche solo coerenza interna. L'alone della viralità crea autorità dal nulla. E una volta che il piedistallo è costruito, il meccanismo descritto sopra si attiva con la stessa potenza che avrebbe per un Nobel: tutto ciò che quella persona dice viene filtrato attraverso la lente della competenza presunta.

Il caso Theranos è forse l'esempio più cinematografico di questo fenomeno portato all'estremo. Elizabeth Holmes, fondatrice dell'azienda, prometteva di rivoluzionare la diagnostica medica con una tecnologia che, semplicemente, non funzionava. Ma l'alone del ``genio visionario della Silicon Valley'' era così potente che investitori navigati, membri di consigli di amministrazione con decenni di esperienza, ex segretari di Stato, continuarono a credere e a investire miliardi. Quando i dubbi emergevano, il ragionamento motivato faceva il suo lavoro: ``Se gente così esperta ci crede, chi sono io per dubitare?''\footnote{Carreyrou, J. (2018). \textit{Bad Blood: Secrets and Lies in a Silicon Valley Startup}. Knopf.}

\section{L'antidoto: ammirare senza accecarsi}

Come ci si libera, allora, dal piedistallo cognitivo? La prima tentazione è il cinismo: non fidarsi di nessuno, trattare ogni autorità con sospetto, assumere che tutti siano incompetenti o in malafede. Ma questo è solo il piedistallo capovolto, un errore uguale e contrario, e altrettanto pericoloso.

La vera soluzione sta in quello che potremmo chiamare \textit{ammirazione granulare}: la capacità di riconoscere e celebrare l'eccellenza specifica di una persona senza estenderla automaticamente a tutto ciò che dice e fa. Einstein era un genio della fisica. Questo non lo rendeva un esperto di relazioni umane (fu, per sua stessa ammissione, un padre assente e un marito difficile). Newton rivoluzionò la meccanica e l'ottica, ma spese anni della sua vita inseguendo il sogno dell'alchimia e cercando codici nascosti nella Bibbia. Nikola Tesla, il genio dell'elettricità, sviluppò negli ultimi anni un legame affettivo con un piccione che lui stesso descriveva in termini quasi mistici.

Questi non sono dettagli gossip. Sono promemoria preziosi del fatto che la genialità non è mai totale, che l'eccellenza in un ambito convive con la normalità, e talvolta con l'errore, in tutti gli altri.

Un secondo strumento è il \textit{principio del secondo parere}. Quando una persona che ammirate si pronuncia su un argomento, specialmente fuori dal suo campo di competenza, cercate attivamente una voce che la pensi diversamente. Non per spirito di contraddizione, ma per fornire al vostro cervello le informazioni che l'effetto alone gli impedisce di cercare spontaneamente. È un po' come quando un buon medico suggerisce di consultare un altro specialista: non perché dubiti di sé stesso, ma perché sa che un secondo sguardo riduce la probabilità di errore.

 Le persone che più meriterebbero la nostra ammirazione incondizionata sono spesso quelle che più resistono a riceverla. Chi conosce davvero la complessità di un campo sa quanto sia facile sbagliare, e quindi dubita, ammette i propri limiti, si corregge. Al contrario, chi sale con entusiasmo sul piedistallo, chi coltiva l'aura dell'infallibilità, chi non ammette mai un errore, è statisticamente più probabile che meriti meno la nostra fiducia.

È un meccanismo che gli psicologi conoscono bene: si chiama effetto Dunning-Kruger.\footnote{Kruger, J., \& Dunning, D. (1999). ``Unskilled and unaware of it: How difficulties in recognizing one's own incompetence lead to inflated self-assessments''. \textit{Journal of Personality and Social Psychology}, 77(6), 1121--1134.} Chi sa poco tende a sopravvalutare le proprie competenze. Chi sa molto tende a sottovalutarle. Risultato: in un mondo che premia la sicurezza e punisce il dubbio, finiamo spesso per mettere sul piedistallo le persone sbagliate.

Richard Feynman, uno dei fisici più brillanti del Novecento, amava ripetere: ``Il primo principio è che non devi ingannare te stesso, e tu sei la persona più facile da ingannare.'' Notate la struttura di quella frase: non sta parlando degli altri. Sta parlando di sé stesso. È l'atteggiamento di chi sa che l'intelligenza non protegge dagli errori. Se mai, li rende più sofisticati e più difficili da individuare.

\section{Il piedistallo vuoto}

Alla fine, il piedistallo cognitivo non è solo un altro bias da aggiungere al catalogo. È, in un certo senso, il meta-bias: quello che ci impedisce di riconoscere gli errori negli altri e, di riflesso, in noi stessi. Quando idealizziamo qualcuno al punto da non vederne le cadute, perdiamo l'occasione di imparare da quelle cadute. E rinunciamo, un passo alla volta, alla cosa più preziosa che abbiamo: la capacità di pensare con la nostra testa.

Ma la consapevolezza è già un primo passo. Sapere che il nostro cervello tende a costruire piedistalli e poi a difenderli con ogni mezzo è, di per sé, una forma di protezione. Non perfetta, certo. Non infallibile. Ma sufficiente a farci fare, ogni tanto, la domanda più importante: ``Sto ammirando questa persona, o sto rinunciando a ragionare?''

Ricordate: anche Newton cercava la pietra filosofale. Anche Pauling sbagliava sulla vitamina C. Anche i giganti, visti da vicino, hanno i piedi d'argilla. E va bene così.\emph{ Perché la vera ammirazione non sta nel credere che qualcuno sia perfetto. Sta nel riconoscere che qualcuno, pur essendo imperfetto come tutti noi, è riuscito a fare qualcosa di straordinario.}

Il piedistallo è vuoto. Lo è sempre stato. E forse è proprio per questo che è così importante continuare a guardare verso l'alto, ma con gli occhi aperti.